\documentclass[11pt,respuestas,a4]{aleph-examen}
% \documentclass[11pt,a4]{aleph-examen}

% -- Paquetes extra
\usepackage{aleph-comandos}
\usepackage{systeme}

% -- Datos
\institucion{Facultad de Ciencias Exactas, Naturales y Ambientales}
\carrera{Catálogo STEM}
\asignatura{Matemática III}
\tema{Repaso -- Criterio 3.2}
\autor{Andrés Merino}
\fecha{Periodo 2026-1}

\logouno[0.16\textwidth]{Logos/logoPUCE_04_ac}
\definecolor{colortext}{HTML}{1957d1}
\definecolor{colordef}{HTML}{1957d1}
\fuente{montserrat}


\begin{document}

\encabezado


\begin{preguntas}

%%%%%%%%%%%%%%%%%%%%%%%%%%%%%%%%%%%%%%%%%%
\item
Encuentre y clasifique todos los puntos críticos de la función
\[
    f(x,y) = x^3 - 3x + y^3 - 3y^2.
\]

\begin{respuesta}
Iniciemos encontrando los puntos críticos. Para ello, calculemos las derivadas parciales de $f$:
\[
    \frac{\partial f}{\partial x}(x,y) = 3x^2 - 3,
    \qquad
    \frac{\partial f}{\partial y}(x,y) = 3y^2 - 6y.
\]
Igualando a cero, obtenemos el sistema:
\[
    \left\{
    \begin{aligned}
        3x^2 - 3 &= 0,\\
        3y^2 - 6y &= 0.
    \end{aligned}
    \right.
\]
Resolviendo, se tiene que $x = \pm 1$ e $y = 0$ o $y = 2$. Con esto, los cuatro puntos críticos son:
\[
    P_1 = (1,0),\quad P_2=(1,2),\quad P_3=(-1,0),\quad P_4=(-1,2).
\]

Apliquemos ahora el criterio de la segunda derivada. Calculemos la matriz Hessiana:
\[
    H_f(x,y)
    =
    \begin{pmatrix}
        \dfrac{\partial^2 f}{\partial x^2} & \dfrac{\partial^2 f}{\partial x\partial y}\\[10pt]
        \dfrac{\partial^2 f}{\partial y\partial x} & \dfrac{\partial^2 f}{\partial y^2}
    \end{pmatrix}
    =
    \begin{pmatrix}
        6x & 0\\
        0 & 6y-6
    \end{pmatrix}.
\]
Evaluando en cada punto crítico y aplicando el criterio de Sylvester:
\begin{itemize}
    \item \textbf{$P_1=(1,0)$:} La Hessiana evaluada en este punto es
    \[
        H_f(1,0) = \begin{pmatrix} 6 & 0 \\ 0 & -6 \end{pmatrix},
    \]
    con esto, $\det(H_f(1,0)) = -36 < 0$, por lo que $(1,0)$ es un punto de silla.

    \item \textbf{$P_2=(1,2)$:} La Hessiana evaluada en este punto es
    \[
        H_f(1,2) = \begin{pmatrix} 6 & 0 \\ 0 & 6 \end{pmatrix},
    \]
    con esto, $\det(H_f(1,2)) = 36 > 0$ y $\dfrac{\partial^2 f}{\partial x^2}(1,2) = 6 > 0$, por lo que la función alcanza un mínimo local en $(1,2)$ y el mínimo es $f(1,2) = -6$.

    \item \textbf{$P_3=(-1,0)$:} La Hessiana evaluada en este punto es
    \[
        H_f(-1,0) = \begin{pmatrix} -6 & 0 \\ 0 & -6 \end{pmatrix},
    \]
    con esto, $\det(H_f(-1,0)) = 36 > 0$ y $\dfrac{\partial^2 f}{\partial x^2}(-1,0) = -6 < 0$, por lo que la función alcanza un máximo local en $(-1,0)$ y el máximo es $f(-1,0) = 2$.

    \item \textbf{$P_4=(-1,2)$:} La Hessiana evaluada en este punto es
    \[
        H_f(-1,2) = \begin{pmatrix} -6 & 0 \\ 0 & 6 \end{pmatrix},
    \]
    con esto, $\det(H_f(-1,2)) = -36 < 0$, por lo que $(-1,2)$ es un punto de silla.\qedhere
\end{itemize}
\end{respuesta}


%%%%%%%%%%%%%%%%%%%%%%%%%%%%%%%%%%%%%%%%%%
\item
Utilice el método de multiplicadores de Lagrange para encontrar los valores extremos de
\[
    f(x,y) = x^2 + 2y^2
\]
sujeta a la restricción $g(x,y) = x + y - 3 = 0$.

\begin{respuesta}
Usemos multiplicadores de Lagrange. Planteemos el Lagrangeano:
\[
    L(x,y,\lambda) = f(x,y) - \lambda\,g(x,y) = x^2 + 2y^2 - \lambda(x+y-3).
\]
Ahora, encontremos los puntos críticos. Para ello, calculemos las derivadas parciales de $L$:
\[
    \frac{\partial L}{\partial x} = 2x - \lambda,
    \qquad
    \frac{\partial L}{\partial y} = 4y - \lambda,
    \qquad
    \frac{\partial L}{\partial \lambda} = -(x+y-3).
\]
Igualando a cero, obtenemos el sistema:
\[
    \left\{
    \begin{aligned}
        2x - \lambda &= 0,\\
        4y - \lambda &= 0,\\
        x + y - 3 &= 0.
    \end{aligned}
    \right.
\]
Resolviendo, se tiene que $x = 2$, $y = 1$ y $\lambda = 4$. Con esto, el único punto crítico es $(2,1)$.

Apliquemos ahora el criterio de la segunda derivada. Calculemos la matriz Hessiana de $L$ con respecto a $(\lambda, x, y)$:
\[
    H_L(\lambda,x,y)
    =
    \begin{pmatrix}
        \dfrac{\partial^2 L}{\partial \lambda^2} & \dfrac{\partial^2 L}{\partial \lambda\,\partial x} & \dfrac{\partial^2 L}{\partial \lambda\,\partial y}\\[10pt]
        \dfrac{\partial^2 L}{\partial x\,\partial \lambda} & \dfrac{\partial^2 L}{\partial x^2} & \dfrac{\partial^2 L}{\partial x\,\partial y}\\[10pt]
        \dfrac{\partial^2 L}{\partial y\,\partial \lambda} & \dfrac{\partial^2 L}{\partial y\,\partial x} & \dfrac{\partial^2 L}{\partial y^2}
    \end{pmatrix}
    =
    \begin{pmatrix}
        0 & -1 & -1\\
        -1 & 2 & 0\\
        -1 & 0 & 4
    \end{pmatrix}.
\]
Evaluando en el punto crítico $(4,2,1)$:
\[
    H_L(4,2,1)
    =
    \begin{pmatrix}
        0 & -1 & -1\\
        -1 & 2 & 0\\
        -1 & 0 & 4
    \end{pmatrix},
\]
con esto, $\det(H_L(4,2,1)) = -6 < 0$, por lo que $f$ alcanza un \textbf{mínimo} en $(2,1)$ y el mínimo es
\[
    f(2,1) = 4 + 2 = 6.\qedhere
\]
\end{respuesta}


%%%%%%%%%%%%%%%%%%%%%%%%%%%%%%%%%%%%%%%%%%
\item
En regresión lineal simple se dispone de $n$ pares de datos $(x_i, y_i)$ y se busca la recta $\hat{y} = ax + b$ que minimiza el \emph{error cuadrático medio}:
\[
    E(a,b)
    =
    \frac{1}{n}\sum_{i=1}^{n}(y_i - ax_i - b)^2.
\]
Para el conjunto de tres puntos $\{(1,2),\,(2,3),\,(3,5)\}$:
\begin{enumerate}
    \item Escriba explícitamente $E(a,b)$ expandida para estos datos.
    \item Encuentre el punto crítico de $E$ resolviendo $\nabla E = {0}$.
    \item Verifique mediante la Hessiana que el punto crítico es un mínimo.
    \item Interprete los coeficientes $a$ y $b$ obtenidos en el contexto del modelo.
\end{enumerate}

\begin{respuesta}

\end{respuesta}

%%%%%%%%%%%%%%%%%%%%%%%%%%%%%%%%%%%%%%%%%%
\item
Una empresa fabrica cajas rectangulares sin tapa con volumen fijo $V_0 = 32\ \text{cm}^3$. El costo del material es de $\$2$ por cm$^2$ para la base y de $\$1$ por cm$^2$ para cada cara lateral. Determine las dimensiones de la caja que minimizan el costo total de material.

\begin{respuesta}
Sean las variables:
\begin{itemize}
    \item $x$: largo de la caja (en cm),
    \item $y$: ancho de la caja (en cm),
    \item $z$: altura de la caja (en cm).
\end{itemize}
El costo total es la suma del costo de la base más el de las cuatro caras laterales:
\[
    f(x,y,z)
    =
    \underbrace{2xy}_{\text{base}}
    + \underbrace{xz + xz}_{\text{front./post.}}
    + \underbrace{yz + yz}_{\text{lat. iz./der.}}
    =
    2xy + 2xz + 2yz.
\]
Como el volumen debe ser de 32 $\text{cm}^3$, se tiene que:
\[
    xyz = 32
\]
Así, el problema de optimización es:
\[
    \left\{
    \begin{aligned}
        &\text{min }  f(x,y,z) = 2xy + 2xz + 2yz,\\
        &\text{sujeto a} \\
        &\qquad xyz - 32 = 0.
    \end{aligned}
    \right.
\]

Usemos multiplicadores de Lagrange. Planteemos el Lagrangeano:
\[
    L(\lambda,x,y,z) = 2xy + 2xz + 2yz - \lambda(xyz - 32).
\]
Ahora, encontremos los puntos críticos. Para ello, calculemos las derivadas parciales de $L$:
\[
    \frac{\partial L}{\partial x} = 2y + 2z - \lambda yz,\quad
    \frac{\partial L}{\partial y} = 2x + 2z - \lambda xz,\quad
    \frac{\partial L}{\partial z} = 2x + 2y - \lambda xy,\quad
    \frac{\partial L}{\partial \lambda} = -(xyz-32).
\]
Igualando a cero, obtenemos el sistema:
\[
    \left\{
    \begin{aligned}
        2y + 2z - \lambda yz &= 0,\\
        2x + 2z - \lambda xz &= 0,\\
        2x + 2y - \lambda xy &= 0,\\
        xyz - 32 &= 0.
    \end{aligned}
    \right.
\]
Resolviendo, se tiene que $x = y = z = 2\sqrt[3]{4}$ y $\lambda = \sqrt[3]{2}$. Con esto, el único punto crítico es $\bigl(2\sqrt[3]{4},\,2\sqrt[3]{4},\,2\sqrt[3]{4}\bigr)$.

Apliquemos ahora el criterio de la segunda derivada. Calculemos la matriz Hessiana de $L$ con respecto a $(\lambda,x,y,z)$:
\[
    H_L(\lambda,x,y,z)
    =
    \begin{pmatrix}
        0 & -yz & -xz & -xy\\[4pt]
        -yz & 0 & 2-\lambda z & 2-\lambda y\\[4pt]
        -xz & 2-\lambda z & 0 & 2-\lambda x\\[4pt]
        -xy & 2-\lambda y & 2-\lambda x & 0
    \end{pmatrix}.
\]
Evaluando en el punto crítico y usando que $\sqrt[3]{2}\cdot 2\sqrt[3]{4} = 4$ y $(2\sqrt[3]{4})^2 = 4\sqrt[3]{16}$:
\[
    H_L\!\left(\sqrt[3]{2},\,2\sqrt[3]{4},\,2\sqrt[3]{4},\,2\sqrt[3]{4}\right)
    =
    \begin{pmatrix}
        0 & -4\sqrt[3]{16} & -4\sqrt[3]{16} & -4\sqrt[3]{16}\\[4pt]
        -4\sqrt[3]{16} & 0 & -2 & -2\\[4pt]
        -4\sqrt[3]{16} & -2 & 0 & -2\\[4pt]
        -4\sqrt[3]{16} & -2 & -2 & 0
    \end{pmatrix}.
\]
Los menores principales relevantes con $m=1$, $n=3$ son:
\[
    \det(H_3) = -4(4\sqrt[3]{16})^2,
    \qquad
    \det(H_4) = -12(4\sqrt[3]{16})^2.
\]
La secuencia del criterio es:
\[
    (-1)^1\det(H_3) = 4(4\sqrt[3]{16})^2 > 0,
    \qquad
    (-1)^1\det(H_4) = 12(4\sqrt[3]{16})^2 > 0.
\]
Como la secuencia consta solo de números positivos, $f$ alcanza un {mínimo} en $\bigl(2\sqrt[3]{4},\,2\sqrt[3]{4},\,2\sqrt[3]{4}\bigr)$ y el costo mínimo es
\[
    f\!\left(2\sqrt[3]{4},\,2\sqrt[3]{4},\,2\sqrt[3]{4}\right)
    =
    6\cdot\bigl(2\sqrt[3]{4}\bigr)^2
    =
    24\sqrt[3]{16}
    \approx \$60{,}48.\qedhere
\]
\end{respuesta}

\end{preguntas}

\end{document}
